% FILE: 01_research_question.tex
% Project: otel-weaver-exp-002-diff-to-residual
% Thesis Role: weaver-diff-to-pi-residual-bridge-experiment

\section{Research Question}
\label{sec:otel-weaver-exp-002-diff-to-residual-rq}

\subsection{Problem Statement}
\label{sec:otel-weaver-exp-002-diff-to-residual-rq-problem}

Process intelligence systems that consume OpenTelemetry telemetry feedstock face a
class of false-positive conformance violation that arises not from actual behavioral
deviation but from schema version misalignment. When an OTel Weaver registry is
updated---for example, to rename an attribute from \texttt{process.pi.activity.name}
to \texttt{process.pi.transition.name} in order to align with Petri net
nomenclature---the conformance court receives telemetry records whose field names no
longer match the court's parser expectations. If no translation layer exists between
the ingestion boundary and the conformance algorithm, the court will compute a reduced
fitness score for those records, interpreting the missing \texttt{process.pi.activity.name}
field as evidence of an unlawful activity omission.

This problem is documented in \texttt{doctrine/registry-diff-is-not-process-drift.md}
(anchored at \texttt{REGISTRY\_DIFF\_NOT\_DRIFT\_001}) as the ``false positive'' hazard:
an organization updates its telemetry schema and the process mining dashboard reports
a sudden compliance failure because the expected activity label has ``disappeared.''
The process is unchanged; only the telemetry label structure drifted. Without an
architectural mechanism to intercept and correct this, the conformance court becomes
polluted by measurement theater---schema noise masquerading as operational risk.

The problem is further sharpened by the addition case: when schema S2 introduces
a new field (\texttt{process.pi.token.color} for Colored Petri Net token tracking),
the translation layer must also decide whether to pass the new field downstream, drop
it with a recorded \texttt{LossReport}, or map it to a corresponding S1 concept. The
current implementation in \texttt{src/bridge.rs} drops \texttt{token\_color} silently
during \texttt{translate()}, which is an open question the experiment has not yet
resolved with documented intent.

\subsection{Old-Regime Assumption Challenged}
\label{sec:otel-weaver-exp-002-diff-to-residual-rq-old-regime}

The old-regime assumption challenged by EXP-002 is: \emph{telemetry schema
compliance is a sufficient proxy for process conformance}. This assumption manifests
in architectures that route OTel validation results directly into compliance
dashboards without a separate process mining layer. Under this assumption, if the
telemetry collector accepts a span without error, the process it describes is assumed
to have executed lawfully.

The doctrine file \texttt{otel-weaver-is-feedstock.md} (anchored at
\texttt{OTEL\_WEAVER\_FEEDSTOCK\_001}) articulates the refutation: schema validation
verifies that instruments are recording within expected structural parameters, not that
the recorded process followed its normative model. A \texttt{dispatch\_payment} span
can pass every OTel Weaver validation rule while describing a payment dispatched
before a credit check was completed---a causal ordering violation that Weaver has no
concept of. The old regime conflates the engineering metric (structural field
presence) with the legal assertion (lawful execution order).

EXP-002 operationalizes this refutation at the specific juncture where schema evolution
makes the conflation most dangerous: a Weaver registry update that renames fields
will cause the old regime's combined schema-plus-conformance check to fire a false
compliance alert on every record emitted under the new schema, until the parser is
manually updated.

\subsection{Hypothesis}
\label{sec:otel-weaver-exp-002-diff-to-residual-rq-hypothesis}

\textbf{Hypothesis:} A \texttt{BridgeRx} translation layer placed at the ingestion
boundary, which maps S2 telemetry feedstock back to S1-compatible structures before
the conformance algorithm fires, will drive the PI Conformance Residual to zero:
\[
    \mathcal{R}_{\Delta} = |f_{S_1}(\mathcal{L}_2, M) - f_{S_2}(\mathcal{L}_2, M)| = 0
\]

This hypothesis claims that the architectural placement of the translation bridge is
sufficient to neutralize schema-evolution noise in the conformance signal, without
requiring any modification to the normative process model $M$ or to the conformance
algorithm itself. The bridge operates purely on the feedstock before it reaches the
court, preserving the court's independence from ingestion concerns.

A secondary hypothesis is that this translation can be implemented using Rust's type
system such that the compiler enforces the categorical separation at compile time:
\texttt{TelemetryFeedstockS2} and \texttt{TelemetryFeedstockS1} are distinct types,
and the \texttt{translate()} method is the only lawful path from one to the other.

\subsection{Success Criteria}
\label{sec:otel-weaver-exp-002-diff-to-residual-rq-success}

Success for EXP-002 requires all of the following, as defined in the broader
checkpoint framework at \texttt{otel-weaver/checkpoints/}:

\begin{enumerate}
    \item \textbf{Compilation}: The \texttt{bridge} library target must compile without
          errors under \texttt{cargo build}.
    \item \textbf{Unit test passage}: \texttt{test\_bridge\_translation\_equivalence}
          must pass, asserting that an S2 record with
          \texttt{transition\_name = "approve\_loan"} is translated to an S1 record
          with \texttt{activity\_name = "approve\_loan"} and that the BLAKE3
          witness hash is preserved through translation.
    \item \textbf{Integration test passage}: The \texttt{BridgeRx} translation path
          must pass within the \texttt{wasm4pm} integration suite referenced in
          \texttt{GGEN\_OTEL\_WEAVER\_PI\_RUNTIME\_001.md}.
    \item \textbf{Documented loss semantics}: The \texttt{token\_color} field drop
          during \texttt{translate()} must be documented as either intentional lossy
          projection with a \texttt{LossReport} receipt, or remediated by a downstream
          mapping.
    \item \textbf{Ontology surface}: A TTL/RDF surface defining the S1/S2 schema types
          and the \texttt{BridgeRx} translation must exist to connect this experiment
          to the OCEL/process-intelligence ontology layer. (\emph{As of the current
          status this criterion is unmet --- see Section~\ref{sec:otel-weaver-exp-002-diff-to-residual-open}.})
\end{enumerate}

Criteria 1 and 2 are partially satisfied by existing evidence. Criteria 3, 4, and 5
remain open, placing the experiment's current status at PARTIAL rather than ALIVE.
